% --- Archon physics formalization source begin ---
% archon:physics
% archon:covers PhyXMiniProblems/problem_phyx_mini_0901.lean
% archon:source-report reports/phyx_mini/problem_phyx_mini_0901.source.json

\chapter{Physics problem phyx\_mini\_0901}
\label{ch:PhyXMiniProblems_problem_phyx_mini_0901}

\paragraph{Problem source.}
\#\# Physical scenario

Two \textbackslash{}( 5.0\textbackslash{},\textbackslash{}mathrm\{g\} \textbackslash{}) point charges on \textbackslash{}( 1.0\textbackslash{},\textbackslash{}mathrm\{m\} \textbackslash{})-long threads repel each other after being charged to \textbackslash{}( +100\textbackslash{},\textbackslash{}mathrm\{nC\} \textbackslash{}), as shown in figure.

\#\# Question

What is the angle \textbackslash{}( 	heta \textbackslash{})?

\#\# Answer choices

- A: A: 3\textasciicircum{}\textbackslash{}circ

- B: B: 9\textasciicircum{}\textbackslash{}circ

- C: C: 4.4\textasciicircum{}\textbackslash{}circ

- D: D: 2\textasciicircum{}\textbackslash{}circ

\#\# Figure caption (auxiliary; use the image as primary evidence)

The image shows a physical setup with two identical objects suspended from a common point, forming an inverted "V" shape. Each object is represented as a red circle with a plus sign in the center, indicating a positive charge. The details are as follows:

- Two identical objects, each with a mass of 5.0 grams and a charge of 100 nanocoulombs (nC).
- Each object is suspended from a point by strings of equal length, 1.0 meter each.
- The angle between each string and the vertical axis is labeled as θ.
- The strings are symmetrically positioned, and a dashed vertical line runs between them, representing the vertical axis.
- The two angles θ are marked on either side of the centerline.

The setup appears to depict a situation where two charged objects are in equilibrium due to electrical and gravitational forces.

\#\# Dataset metadata

Electrostatics; Physical Model Grounding Reasoning, Multi-Formula Reasoning

\paragraph{Recorded answer/context.}
C: C: 4.4\textasciicircum{}\textbackslash{}circ

\paragraph{Figure/image path.}
/root/proposal\_for\_physic/science-mango/phyx\_mini\_run/phyx\_data/test\_image/901.png

\paragraph{Formalization target.}
create a compiling Lean file with sorry bodies at `PhyXMiniProblems/problem\_phyx\_mini\_0901.lean`. The Lean declarations must preserve the physical quantities, dimensions or dimensional roles, figure labels, governing-law hypotheses, and final relation expressed by this problem.
Use Mathlib/Physlib names found through LeanExplore where available. If a physics API is missing, introduce faithful local abstractions rather than scalar placeholder aliases.

\begin{theorem}[Physics formalization target]
\label{thm:physics:phyx_mini_0901:target}
\lean{PhyXMiniProblems.ProblemPhyXMini0901.problem_phyx_mini_0901}
\uses{def:physics:phyx-mini-0901:phyxminiproblems-problemphyxmini0901-chargedpendulumsetup, def:physics:phyx-mini-0901:phyxminiproblems-problemphyxmini0901-matcheschargedpendulumscenario, def:physics:phyx-mini-0901:phyxminiproblems-problemphyxmini0901-matchessuppliedfigure, def:physics:phyx-mini-0901:phyxminiproblems-problemphyxmini0901-hasproblemphysicaldata, def:physics:phyx-mini-0901:phyxminiproblems-problemphyxmini0901-hastextbookconstantcalibrations, def:physics:phyx-mini-0901:phyxminiproblems-problemphyxmini0901-hasphysicalchargedpendulumparameters, def:physics:phyx-mini-0901:phyxminiproblems-problemphyxmini0901-satisfiessymmetricpendulumgeometry, def:physics:phyx-mini-0901:phyxminiproblems-problemphyxmini0901-satisfieschargedpendulumequilibriumlaws, def:physics:phyx-mini-0901:phyxminiproblems-problemphyxmini0901-ismatchinganswer}
The assigned autoformalize agent should translate this physics problem into Lean declarations in the covered file, with theorem and lemma proof bodies written as `by sorry`.
\end{theorem}
\begin{proof}
This is an autoformalization task, not a proof task. Produce statements that can later be proved by the physics prover without weakening the physical model.
\end{proof}

% --- Archon declaration topology begin ---
\section{Declaration topology}
\begin{definition}[acceleration Dimension]
  \label{def:physics:phyx-mini-0901:phyxminiproblems-problemphyxmini0901-accelerationdimension}
  \lean{PhyXMiniProblems.ProblemPhyXMini0901.accelerationDimension}
  The physical dimension L T⁻² of acceleration
\end{definition}
\begin{proof}
  This is the declared model definition of acceleration Dimension.
\end{proof}
\begin{definition}[force Dimension]
  \label{def:physics:phyx-mini-0901:phyxminiproblems-problemphyxmini0901-forcedimension}
  \lean{PhyXMiniProblems.ProblemPhyXMini0901.forceDimension}
  The physical dimension M L T⁻² of force
\end{definition}
\begin{proof}
  This is the declared model definition of force Dimension.
\end{proof}
\begin{definition}[coulomb Constant Dimension]
  \label{def:physics:phyx-mini-0901:phyxminiproblems-problemphyxmini0901-coulombconstantdimension}
  \lean{PhyXMiniProblems.ProblemPhyXMini0901.coulombConstantDimension}
  \uses{def:physics:phyx-mini-0901:phyxminiproblems-problemphyxmini0901-forcedimension}
  The dimension force * length² / charge² of Coulomb's constant
\end{definition}
\begin{proof}
  This is the declared model definition of coulomb Constant Dimension.
\end{proof}
\begin{definition}[Mass Quantity]
  \label{def:physics:phyx-mini-0901:phyxminiproblems-problemphyxmini0901-massquantity}
  \lean{PhyXMiniProblems.ProblemPhyXMini0901.MassQuantity}
  A nonnegative mass as a unit-independent physical quantity
\end{definition}
\begin{proof}
  This is the declared model definition of Mass Quantity.
\end{proof}
\begin{definition}[Length Quantity]
  \label{def:physics:phyx-mini-0901:phyxminiproblems-problemphyxmini0901-lengthquantity}
  \lean{PhyXMiniProblems.ProblemPhyXMini0901.LengthQuantity}
  A nonnegative length as a unit-independent physical quantity
\end{definition}
\begin{proof}
  This is the declared model definition of Length Quantity.
\end{proof}
\begin{definition}[Signed Charge Quantity]
  \label{def:physics:phyx-mini-0901:phyxminiproblems-problemphyxmini0901-signedchargequantity}
  \lean{PhyXMiniProblems.ProblemPhyXMini0901.SignedChargeQuantity}
  A signed electric charge as a unit-independent physical quantity
\end{definition}
\begin{proof}
  This is the declared model definition of Signed Charge Quantity.
\end{proof}
\begin{definition}[Acceleration Magnitude Quantity]
  \label{def:physics:phyx-mini-0901:phyxminiproblems-problemphyxmini0901-accelerationmagnitudequantity}
  \lean{PhyXMiniProblems.ProblemPhyXMini0901.AccelerationMagnitudeQuantity}
  \uses{def:physics:phyx-mini-0901:phyxminiproblems-problemphyxmini0901-accelerationdimension}
  A nonnegative acceleration magnitude
\end{definition}
\begin{proof}
  This is the declared model definition of Acceleration Magnitude Quantity.
\end{proof}
\begin{definition}[Force Magnitude Quantity]
  \label{def:physics:phyx-mini-0901:phyxminiproblems-problemphyxmini0901-forcemagnitudequantity}
  \lean{PhyXMiniProblems.ProblemPhyXMini0901.ForceMagnitudeQuantity}
  \uses{def:physics:phyx-mini-0901:phyxminiproblems-problemphyxmini0901-forcedimension}
  A nonnegative force magnitude
\end{definition}
\begin{proof}
  This is the declared model definition of Force Magnitude Quantity.
\end{proof}
\begin{definition}[Coulomb Constant Quantity]
  \label{def:physics:phyx-mini-0901:phyxminiproblems-problemphyxmini0901-coulombconstantquantity}
  \lean{PhyXMiniProblems.ProblemPhyXMini0901.CoulombConstantQuantity}
  \uses{def:physics:phyx-mini-0901:phyxminiproblems-problemphyxmini0901-coulombconstantdimension}
  A nonnegative physical value of Coulomb's constant
\end{definition}
\begin{proof}
  This is the declared model definition of Coulomb Constant Quantity.
\end{proof}
\begin{definition}[mass In Kilograms]
  \label{def:physics:phyx-mini-0901:phyxminiproblems-problemphyxmini0901-massinkilograms}
  \lean{PhyXMiniProblems.ProblemPhyXMini0901.massInKilograms}
  \uses{def:physics:phyx-mini-0901:phyxminiproblems-problemphyxmini0901-massquantity}
  Read a physical mass in kilograms
\end{definition}
\begin{proof}
  This is the declared model definition of mass In Kilograms.
\end{proof}
\begin{definition}[length In Meters]
  \label{def:physics:phyx-mini-0901:phyxminiproblems-problemphyxmini0901-lengthinmeters}
  \lean{PhyXMiniProblems.ProblemPhyXMini0901.lengthInMeters}
  \uses{def:physics:phyx-mini-0901:phyxminiproblems-problemphyxmini0901-lengthquantity}
  Read a physical length in metres
\end{definition}
\begin{proof}
  This is the declared model definition of length In Meters.
\end{proof}
\begin{definition}[charge In Coulombs]
  \label{def:physics:phyx-mini-0901:phyxminiproblems-problemphyxmini0901-chargeincoulombs}
  \lean{PhyXMiniProblems.ProblemPhyXMini0901.chargeInCoulombs}
  \uses{def:physics:phyx-mini-0901:phyxminiproblems-problemphyxmini0901-signedchargequantity}
  Read a signed physical charge in coulombs
\end{definition}
\begin{proof}
  This is the declared model definition of charge In Coulombs.
\end{proof}
\begin{definition}[acceleration In Meters Per Second Squared]
  \label{def:physics:phyx-mini-0901:phyxminiproblems-problemphyxmini0901-accelerationinmeterspersecondsquared}
  \lean{PhyXMiniProblems.ProblemPhyXMini0901.accelerationInMetersPerSecondSquared}
  \uses{def:physics:phyx-mini-0901:phyxminiproblems-problemphyxmini0901-accelerationmagnitudequantity}
  Read an acceleration magnitude in metres per second squared
\end{definition}
\begin{proof}
  This is the declared model definition of acceleration In Meters Per Second Squared.
\end{proof}
\begin{definition}[force In Newtons]
  \label{def:physics:phyx-mini-0901:phyxminiproblems-problemphyxmini0901-forceinnewtons}
  \lean{PhyXMiniProblems.ProblemPhyXMini0901.forceInNewtons}
  \uses{def:physics:phyx-mini-0901:phyxminiproblems-problemphyxmini0901-forcemagnitudequantity}
  Read a force magnitude in newtons
\end{definition}
\begin{proof}
  This is the declared model definition of force In Newtons.
\end{proof}
\begin{definition}[coulomb Constant In Newton Meters Squared Per Coulomb Squared]
  \label{def:physics:phyx-mini-0901:phyxminiproblems-problemphyxmini0901-coulombconstantinnewtonmeterssquaredpercoulombsquared}
  \lean{PhyXMiniProblems.ProblemPhyXMini0901.coulombConstantInNewtonMetersSquaredPerCoulombSquared}
  \uses{def:physics:phyx-mini-0901:phyxminiproblems-problemphyxmini0901-coulombconstantquantity}
  Read Coulomb's constant in N m² / C²
\end{definition}
\begin{proof}
  This is the declared model definition of coulomb Constant In Newton Meters Squared Per Coulomb Squared.
\end{proof}
\begin{definition}[angle In Degrees]
  \label{def:physics:phyx-mini-0901:phyxminiproblems-problemphyxmini0901-angleindegrees}
  \lean{PhyXMiniProblems.ProblemPhyXMini0901.angleInDegrees}
  Read the principal representative of an angle in degrees
\end{definition}
\begin{proof}
  This is the declared model definition of angle In Degrees.
\end{proof}
\begin{definition}[Bob]
  \label{def:physics:phyx-mini-0901:phyxminiproblems-problemphyxmini0901-bob}
  \lean{PhyXMiniProblems.ProblemPhyXMini0901.Bob}
  The left and right suspended charged bodies
\end{definition}
\begin{proof}
  This is the declared model definition of Bob.
\end{proof}
\begin{definition}[Charged Body Model]
  \label{def:physics:phyx-mini-0901:phyxminiproblems-problemphyxmini0901-chargedbodymodel}
  \lean{PhyXMiniProblems.ProblemPhyXMini0901.ChargedBodyModel}
  Idealization of the two charged bodies used by Coulomb's law
\end{definition}
\begin{proof}
  This is the declared model definition of Charged Body Model.
\end{proof}
\begin{definition}[Thread Model]
  \label{def:physics:phyx-mini-0901:phyxminiproblems-problemphyxmini0901-threadmodel}
  \lean{PhyXMiniProblems.ProblemPhyXMini0901.ThreadModel}
  Idealization of each supporting thread
\end{definition}
\begin{proof}
  This is the declared model definition of Thread Model.
\end{proof}
\begin{definition}[Mechanical State]
  \label{def:physics:phyx-mini-0901:phyxminiproblems-problemphyxmini0901-mechanicalstate}
  \lean{PhyXMiniProblems.ProblemPhyXMini0901.MechanicalState}
  Mechanical state shown after the charges have separated
\end{definition}
\begin{proof}
  This is the declared model definition of Mechanical State.
\end{proof}
\begin{definition}[Printed Charge Sign]
  \label{def:physics:phyx-mini-0901:phyxminiproblems-problemphyxmini0901-printedchargesign}
  \lean{PhyXMiniProblems.ProblemPhyXMini0901.PrintedChargeSign}
  Signs printed inside the two red charge markers
\end{definition}
\begin{proof}
  This is the declared model definition of Printed Charge Sign.
\end{proof}
\begin{definition}[Figure Object]
  \label{def:physics:phyx-mini-0901:phyxminiproblems-problemphyxmini0901-figureobject}
  \lean{PhyXMiniProblems.ProblemPhyXMini0901.FigureObject}
  Physical or geometric objects visibly present in the supplied image
\end{definition}
\begin{proof}
  This is the declared model definition of Figure Object.
\end{proof}
\begin{definition}[Figure Label]
  \label{def:physics:phyx-mini-0901:phyxminiproblems-problemphyxmini0901-figurelabel}
  \lean{PhyXMiniProblems.ProblemPhyXMini0901.FigureLabel}
  Literal label kinds displayed in image 901.png
\end{definition}
\begin{proof}
  This is the declared model definition of Figure Label.
\end{proof}
\begin{definition}[Charged Pendulum Figure]
  \label{def:physics:phyx-mini-0901:phyxminiproblems-problemphyxmini0901-chargedpendulumfigure}
  \lean{PhyXMiniProblems.ProblemPhyXMini0901.ChargedPendulumFigure}
  \uses{def:physics:phyx-mini-0901:phyxminiproblems-problemphyxmini0901-bob, def:physics:phyx-mini-0901:phyxminiproblems-problemphyxmini0901-printedchargesign, def:physics:phyx-mini-0901:phyxminiproblems-problemphyxmini0901-figureobject, def:physics:phyx-mini-0901:phyxminiproblems-problemphyxmini0901-figurelabel}
  Presentation data transcribed from the primary image
\end{definition}
\begin{proof}
  This is the declared model definition of Charged Pendulum Figure.
\end{proof}
\begin{definition}[Charged Pendulum Setup]
  \label{def:physics:phyx-mini-0901:phyxminiproblems-problemphyxmini0901-chargedpendulumsetup}
  \lean{PhyXMiniProblems.ProblemPhyXMini0901.ChargedPendulumSetup}
  \uses{def:physics:phyx-mini-0901:phyxminiproblems-problemphyxmini0901-massquantity, def:physics:phyx-mini-0901:phyxminiproblems-problemphyxmini0901-lengthquantity, def:physics:phyx-mini-0901:phyxminiproblems-problemphyxmini0901-signedchargequantity, def:physics:phyx-mini-0901:phyxminiproblems-problemphyxmini0901-accelerationmagnitudequantity, def:physics:phyx-mini-0901:phyxminiproblems-problemphyxmini0901-forcemagnitudequantity, def:physics:phyx-mini-0901:phyxminiproblems-problemphyxmini0901-coulombconstantquantity, def:physics:phyx-mini-0901:phyxminiproblems-problemphyxmini0901-bob, def:physics:phyx-mini-0901:phyxminiproblems-problemphyxmini0901-chargedbodymodel, def:physics:phyx-mini-0901:phyxminiproblems-problemphyxmini0901-threadmodel, def:physics:phyx-mini-0901:phyxminiproblems-problemphyxmini0901-mechanicalstate, def:physics:phyx-mini-0901:phyxminiproblems-problemphyxmini0901-chargedpendulumfigure}
  Independent physical quantities and observables of the two-pendulum system. No force, separation, or angle is defined from the recorded answer
\end{definition}
\begin{proof}
  This is the declared model definition of Charged Pendulum Setup.
\end{proof}
\begin{definition}[Matches Charged Pendulum Scenario]
  \label{def:physics:phyx-mini-0901:phyxminiproblems-problemphyxmini0901-matcheschargedpendulumscenario}
  \lean{PhyXMiniProblems.ProblemPhyXMini0901.MatchesChargedPendulumScenario}
  \uses{def:physics:phyx-mini-0901:phyxminiproblems-problemphyxmini0901-chargedpendulumsetup}
  Qualitative assumptions stated or implied by the physical scenario
\end{definition}
\begin{proof}
  This is the declared model definition of Matches Charged Pendulum Scenario.
\end{proof}
\begin{definition}[Matches Supplied Figure]
  \label{def:physics:phyx-mini-0901:phyxminiproblems-problemphyxmini0901-matchessuppliedfigure}
  \lean{PhyXMiniProblems.ProblemPhyXMini0901.MatchesSuppliedFigure}
  \uses{def:physics:phyx-mini-0901:phyxminiproblems-problemphyxmini0901-chargedpendulumsetup}
  Exact qualitative and literal information visible in image 901.png
\end{definition}
\begin{proof}
  This is the declared model definition of Matches Supplied Figure.
\end{proof}
\begin{definition}[Has Problem Physical Data]
  \label{def:physics:phyx-mini-0901:phyxminiproblems-problemphyxmini0901-hasproblemphysicaldata}
  \lean{PhyXMiniProblems.ProblemPhyXMini0901.HasProblemPhysicalData}
  \uses{def:physics:phyx-mini-0901:phyxminiproblems-problemphyxmini0901-massinkilograms, def:physics:phyx-mini-0901:phyxminiproblems-problemphyxmini0901-lengthinmeters, def:physics:phyx-mini-0901:phyxminiproblems-problemphyxmini0901-chargeincoulombs, def:physics:phyx-mini-0901:phyxminiproblems-problemphyxmini0901-chargedpendulumsetup}
  Dimensionful physical values stated in the problem. The signed charge readout also records that both charges are positive. No angle or force value occurs in this interface
\end{definition}
\begin{proof}
  This is the declared model definition of Has Problem Physical Data.
\end{proof}
\begin{definition}[Has Textbook Constant Calibrations]
  \label{def:physics:phyx-mini-0901:phyxminiproblems-problemphyxmini0901-hastextbookconstantcalibrations}
  \lean{PhyXMiniProblems.ProblemPhyXMini0901.HasTextbookConstantCalibrations}
  \uses{def:physics:phyx-mini-0901:phyxminiproblems-problemphyxmini0901-accelerationinmeterspersecondsquared, def:physics:phyx-mini-0901:phyxminiproblems-problemphyxmini0901-coulombconstantinnewtonmeterssquaredpercoulombsquared, def:physics:phyx-mini-0901:phyxminiproblems-problemphyxmini0901-chargedpendulumsetup}
  Standard textbook calibrations needed for the numerical calculation
\end{definition}
\begin{proof}
  This is the declared model definition of Has Textbook Constant Calibrations.
\end{proof}
\begin{definition}[Has Physical Charged Pendulum Parameters]
  \label{def:physics:phyx-mini-0901:phyxminiproblems-problemphyxmini0901-hasphysicalchargedpendulumparameters}
  \lean{PhyXMiniProblems.ProblemPhyXMini0901.HasPhysicalChargedPendulumParameters}
  \uses{def:physics:phyx-mini-0901:phyxminiproblems-problemphyxmini0901-massinkilograms, def:physics:phyx-mini-0901:phyxminiproblems-problemphyxmini0901-lengthinmeters, def:physics:phyx-mini-0901:phyxminiproblems-problemphyxmini0901-chargeincoulombs, def:physics:phyx-mini-0901:phyxminiproblems-problemphyxmini0901-accelerationinmeterspersecondsquared, def:physics:phyx-mini-0901:phyxminiproblems-problemphyxmini0901-forceinnewtons, def:physics:phyx-mini-0901:phyxminiproblems-problemphyxmini0901-coulombconstantinnewtonmeterssquaredpercoulombsquared, def:physics:phyx-mini-0901:phyxminiproblems-problemphyxmini0901-chargedpendulumsetup}
  Positivity and the acute angular branch depicted in the figure
\end{definition}
\begin{proof}
  This is the declared model definition of Has Physical Charged Pendulum Parameters.
\end{proof}
\begin{definition}[Satisfies Symmetric Pendulum Geometry]
  \label{def:physics:phyx-mini-0901:phyxminiproblems-problemphyxmini0901-satisfiessymmetricpendulumgeometry}
  \lean{PhyXMiniProblems.ProblemPhyXMini0901.SatisfiesSymmetricPendulumGeometry}
  \uses{def:physics:phyx-mini-0901:phyxminiproblems-problemphyxmini0901-lengthinmeters, def:physics:phyx-mini-0901:phyxminiproblems-problemphyxmini0901-chargedpendulumsetup}
  Mirror-symmetric common-pivot geometry. Each bob's horizontal displacement is L sin(theta), so the center-to-center separation is 2 L sin(theta)
\end{definition}
\begin{proof}
  This is the declared model definition of Satisfies Symmetric Pendulum Geometry.
\end{proof}
\begin{definition}[Satisfies Charged Pendulum Equilibrium Laws]
  \label{def:physics:phyx-mini-0901:phyxminiproblems-problemphyxmini0901-satisfieschargedpendulumequilibriumlaws}
  \lean{PhyXMiniProblems.ProblemPhyXMini0901.SatisfiesChargedPendulumEquilibriumLaws}
  \uses{def:physics:phyx-mini-0901:phyxminiproblems-problemphyxmini0901-massinkilograms, def:physics:phyx-mini-0901:phyxminiproblems-problemphyxmini0901-lengthinmeters, def:physics:phyx-mini-0901:phyxminiproblems-problemphyxmini0901-chargeincoulombs, def:physics:phyx-mini-0901:phyxminiproblems-problemphyxmini0901-accelerationinmeterspersecondsquared, def:physics:phyx-mini-0901:phyxminiproblems-problemphyxmini0901-forceinnewtons, def:physics:phyx-mini-0901:phyxminiproblems-problemphyxmini0901-coulombconstantinnewtonmeterssquaredpercoulombsquared, def:physics:phyx-mini-0901:phyxminiproblems-problemphyxmini0901-chargedpendulumsetup}
  Coulomb's inverse-square law, weight m g, and horizontal and vertical component force balance. These are coherent-SI readout equations and contain neither the target angle nor an answer-choice label
\end{definition}
\begin{proof}
  This is the declared model definition of Satisfies Charged Pendulum Equilibrium Laws.
\end{proof}
\begin{definition}[Answer Choice]
  \label{def:physics:phyx-mini-0901:phyxminiproblems-problemphyxmini0901-answerchoice}
  \lean{PhyXMiniProblems.ProblemPhyXMini0901.AnswerChoice}
  Multiple-choice labels in the order displayed by the source
\end{definition}
\begin{proof}
  This is the declared model definition of Answer Choice.
\end{proof}
\begin{definition}[displayed Angle In Degrees]
  \label{def:physics:phyx-mini-0901:phyxminiproblems-problemphyxmini0901-answerchoice-displayedangleindegrees}
  \lean{PhyXMiniProblems.ProblemPhyXMini0901.AnswerChoice.displayedAngleInDegrees}
  \uses{def:physics:phyx-mini-0901:phyxminiproblems-problemphyxmini0901-answerchoice}
  Numerical angle printed beside an answer choice, in degrees
\end{definition}
\begin{proof}
  This is the declared model definition of displayed Angle In Degrees.
\end{proof}
\begin{definition}[unit In Last Displayed Place In Degrees]
  \label{def:physics:phyx-mini-0901:phyxminiproblems-problemphyxmini0901-answerchoice-unitinlastdisplayedplaceindegrees}
  \lean{PhyXMiniProblems.ProblemPhyXMini0901.AnswerChoice.unitInLastDisplayedPlaceInDegrees}
  \uses{def:physics:phyx-mini-0901:phyxminiproblems-problemphyxmini0901-answerchoice}
  Precision of the final printed digit of each displayed choice
\end{definition}
\begin{proof}
  This is the declared model definition of unit In Last Displayed Place In Degrees.
\end{proof}
\begin{definition}[Matches Displayed Precision]
  \label{def:physics:phyx-mini-0901:phyxminiproblems-problemphyxmini0901-matchesdisplayedprecision}
  \lean{PhyXMiniProblems.ProblemPhyXMini0901.MatchesDisplayedPrecision}
  \uses{def:physics:phyx-mini-0901:phyxminiproblems-problemphyxmini0901-answerchoice}
  An actual angle agrees with a displayed choice under round-to-nearest
\end{definition}
\begin{proof}
  This is the declared model definition of Matches Displayed Precision.
\end{proof}
\begin{definition}[Is Matching Answer]
  \label{def:physics:phyx-mini-0901:phyxminiproblems-problemphyxmini0901-ismatchinganswer}
  \lean{PhyXMiniProblems.ProblemPhyXMini0901.IsMatchingAnswer}
  \uses{def:physics:phyx-mini-0901:phyxminiproblems-problemphyxmini0901-angleindegrees, def:physics:phyx-mini-0901:phyxminiproblems-problemphyxmini0901-chargedpendulumsetup, def:physics:phyx-mini-0901:phyxminiproblems-problemphyxmini0901-answerchoice, def:physics:phyx-mini-0901:phyxminiproblems-problemphyxmini0901-matchesdisplayedprecision}
  A choice matches the independently modeled physical equilibrium angle
\end{definition}
\begin{proof}
  This is the declared model definition of Is Matching Answer.
\end{proof}
\begin{definition}[recorded Dataset Answer]
  \label{def:physics:phyx-mini-0901:phyxminiproblems-problemphyxmini0901-recordeddatasetanswer}
  \lean{PhyXMiniProblems.ProblemPhyXMini0901.recordedDatasetAnswer}
  \uses{def:physics:phyx-mini-0901:phyxminiproblems-problemphyxmini0901-answerchoice}
  Answer label recorded in the source dataset; metadata only
\end{definition}
\begin{proof}
  This is the declared model definition of recorded Dataset Answer.
\end{proof}
% --- Archon declaration topology end ---
% --- Archon physics formalization source end ---
